% SPDX-License-Identifier: Apache-2.0
\documentclass[11pt]{article}

\usepackage[margin=1in]{geometry}
\usepackage[T1]{fontenc}
\usepackage[utf8]{inputenc}
\usepackage{array}
\usepackage{booktabs}
\usepackage{longtable}
\usepackage{amsmath}
\usepackage{amssymb}
\usepackage{xcolor}
\usepackage[hidelinks]{hyperref}

\newcommand{\protocol}{Telara\~na}
\newcommand{\repo}{\texttt{fx-telarana}}
\newcommand{\commithead}{\texttt{8de5959}}
\newcommand{\tokenledger}{\texttt{PENDING\_OPENAI\_USAGE\_EXPORT}}

\title{\protocol: A Cross-Chain Onchain Forex Credit Hub}
\author{
  GPT.5.5 Extra High, Main Author\thanks{Prepared as a Codex-authored technical whitepaper for the \repo{} repository at the request of the maintainer. Exact model-token accounting is not available from the repository or local shell runtime; see Appendix~\ref{appendix:token-ledger}.}
  \and Tom\'as Cordero
}
\date{May 15, 2026 \\ Draft v0.1}

\begin{document}
\maketitle

\begin{abstract}
\protocol{} is a cross-chain onchain forex credit hub. The protocol lets users
enter from supported chains with USDC or EURC where Circle supports the route,
move USDC between \protocol{} hubs through Circle Gateway, route into hub FX
markets, and borrow, lend, or prepare spot-FX requests against supported
fiat-stablecoin pairs. Morpho Blue provides the isolated lending substrate; a
single oracle read surface,
\texttt{IFxOracle}, combines Pyth and RedStone feeds with stale-price,
confidence, and deviation checks; Uniswap v4 hooks prepare spot-FX execution,
oracle-aware spread control, and liquidity rehypothecation; CCTP is used only
for Circle-supported USDC/EURC spoke entry; Circle Gateway is USDC-only in the
current design; Hyperlane and approved issuer-specific routes handle other
stablecoin transport and cross-chain intent messages.

The objective is not to build a generic bridge. The objective is to make many
chains feel like one FX credit protocol while keeping the canonical risk engine
on a hub. \protocol{} accepts only registered assets, routes, markets, oracles,
and execution surfaces. The current smart-contract handoff scope has been
completed and verified with local unit tests, mainnet-fork tests, SDK tests,
Gateway edge-case tests, and Uniswap v4 hook invariants.
\end{abstract}

\section*{Status and Disclaimer}

This paper describes a protocol design and the smart-contract implementation
state represented by protocol verification commit \commithead{}. It is not investment,
legal, tax, accounting, regulatory, or security advice. The current protocol
handoff is not a substitute for an external audit, formal verification, live
deployment review, or operational monitoring. Any production deployment that
moves user funds requires independent review, issuer-route verification,
oracle-feed verification, privileged-key management review, and incident
response procedures.

\section{Mission}

The mission of \protocol{} is to create a production-grade, onchain FX credit
primitive:

\begin{quote}
Deposit dollar liquidity, borrow local-currency stablecoins, swap between
currency rails, and manage collateralized credit with institutional-grade risk
controls.
\end{quote}

\protocol{} focuses on three user-visible capabilities:

\begin{enumerate}
  \item Deposit USDC or supported stablecoins as collateral.
  \item Borrow foreign-currency stablecoins such as AUDF, JPYC, MXNB, KRW1, or
        ZCHF where the asset, oracle route, liquidity route, and market are
        approved.
  \item Execute or prepare FX swaps through Uniswap v4 spot-FX surfaces while
        keeping risk accounting anchored to the hub.
\end{enumerate}

The protocol's core claim is simple: cross-chain FX credit should not be
implemented as uncontrolled token wrapping. It should be implemented as
registered intent, registered transport, registered markets, and registered
risk.

\section{System Overview}

\protocol{} is organized as a hub-and-spoke system. Avalanche is the first
mainnet hub target because the current listed basket has unusually strong
coverage there today; it is not the protocol's identity or long-term boundary.
Fuji and Arc Testnet are testnet hubs. Other chains can act as hubs, spokes, or
intent origins depending on asset availability, route safety, and market depth.

\begin{center}
\begin{tabular}{p{0.24\linewidth}p{0.68\linewidth}}
\toprule
Component & Role \\
\midrule
\texttt{FxMarketRegistry} & App-facing surface over Morpho Blue markets:
supply, withdraw, supply collateral, withdraw collateral, borrow, and repay. \\
\texttt{FxOracle} & Pyth-primary and RedStone-secondary oracle with stale,
confidence, and deviation checks. \\
\texttt{MorphoOracleAdapter} & Adapter from \texttt{IFxOracle} into the
Morpho Blue \texttt{IOracle.price()} shape. \\
\texttt{FxSpoke} & CCTP spoke entry for Circle-supported USDC/EURC routes,
with explicit beneficiary and stranded-deposit recovery. \\
\texttt{FxHubMessageReceiver} & CCTP hub receiver that validates inbound
message state and executes hub calldata or marks deposits stranded. \\
\texttt{FxSpokeIntentRouter} & Hyperlane intent sender for route, market, and
action metadata. \\
\texttt{FxHyperlaneHubReceiver} & Hyperlane hub receiver that accepts only
trusted origin spokes, assets, routes, and nonces. \\
\texttt{TelaranaGatewayHubHook} & Circle Gateway destination-hub wrapper for
USDC hub-to-hub mint and optional future spot-FX request emission. \\
\texttt{FxSwapHook} & Uniswap v4 hook for oracle-anchored spot FX, dynamic
spread, truncated observations, and Morpho rehypothecation preparation. \\
\bottomrule
\end{tabular}
\end{center}

\section{Asset and Route Model}

\subsection{Current Basket}

The target basket is:

\begin{center}
\begin{tabular}{llll}
\toprule
Asset & Role & Initial hub status in design & Notes \\
\midrule
USDC & Settlement and collateral anchor & Supported & Circle rail asset. \\
AUDF & AUD stablecoin & Initial hub target & Six decimals in reference data. \\
JPYC & JPY stablecoin & Initial hub target & Mainnet decimals differ from some test deployments. \\
MXNB & MXN stablecoin & Initial hub target & Six decimals in reference data. \\
KRW1 & KRW stablecoin & Initial hub target & Decimals must be verified before mocks. \\
ZCHF & CHF stablecoin & Initial hub target & CCIP/issuer route considerations. \\
EURC & EUR stablecoin & Disabled/TBD in basket scope & CCTP-only where Circle supports it. \\
\bottomrule
\end{tabular}
\end{center}

BRLA and PHPC are excluded from the current active mainnet basket. QCAD and
ZARU remain blocked by issuer, chain, or token-identity constraints.

\subsection{Transport Rules}

\protocol{} distinguishes asset transport from asset acceptance:

\begin{itemize}
  \item Circle Gateway is USDC-only in the current design.
  \item CCTP is used only for Circle-supported USDC/EURC movement.
  \item Hyperlane and approved issuer-specific routes handle other stablecoin
        transport and intent messages.
  \item No permissionless wrapped asset is automatically accepted as
        collateral.
\end{itemize}

For every asset route, the hub must know:

\[
  \text{asset identity, origin, route id, route token, oracle route, market id, risk parameters.}
\]

If any of these fields is absent or disabled, the hub rejects the action.

\section{Lending Architecture}

\protocol{} uses Morpho Blue as a lending substrate. Each FX pair is deployed
as isolated markets, not as a pooled cross-margin system. For a pair such as
USDC/JPYC, the protocol can deploy:

\begin{center}
\begin{tabular}{llll}
\toprule
Market & Loan asset & Collateral asset & Purpose \\
\midrule
M1 & JPYC & USDC & Borrow local currency against dollar collateral. \\
M2 & USDC & JPYC & Borrow dollars against local-currency collateral. \\
\bottomrule
\end{tabular}
\end{center}

The registry keeps app logic away from raw Morpho calls. It also gives
operations a per-pool live switch. A global pause stops entry-side actions, and
the per-pool switch isolates pair-specific incidents without shutting down the
whole protocol.

\subsection{Borrowing and Repayment}

Borrowing is routed through \texttt{FxMarketRegistry}. The registry enforces
the selected pool id, live status, and caller authorization. Repayment uses the
same registry surface, keeping integrations on one app-facing contract.

\subsection{Collateral Receipts}

\texttt{FxReceipt} prepares ERC-4626-style receipt accounting for supported
assets. Receipt surfaces are useful for frontend accounting, indexer state, and
future composability, but they do not replace Morpho's underlying accounting.

\section{Oracle and FX Risk Engine}

Foreign-exchange pairs are less volatile than many crypto assets, but the
protocol still treats every price as adversarial until verified. The risk
engine uses a layered defense:

\begin{enumerate}
  \item \textbf{Single read path:} contracts read price data only through
        \texttt{IFxOracle}.
  \item \textbf{Primary and secondary feeds:} Pyth is the primary source and
        RedStone is the secondary verification route.
  \item \textbf{Staleness checks:} old data is rejected.
  \item \textbf{Confidence checks:} wide confidence intervals are rejected.
  \item \textbf{Deviation checks:} Pyth and RedStone must agree within a
        configured bound for verified reads.
  \item \textbf{Per-market isolation:} market failures are contained by
        isolated Morpho markets and per-pool live switches.
  \item \textbf{Liquidation paths:} positions remain liquidatable under
        Morpho's isolated-market rules.
\end{enumerate}

This is the core design rule:

\begin{quote}
Transport may be permissionless. Collateral acceptance is never permissionless.
\end{quote}

\section{Cross-Chain Entry}

\subsection{CCTP for USDC/EURC}

\texttt{FxSpoke} is the thin CCTP spoke contract. It accepts only Circle assets
that are explicitly enabled for the route. It takes an explicit
\texttt{beneficiary}; it does not infer ownership from \texttt{msg.sender}.
This is required for both public entry and future Ghost Mode accounts.

\texttt{FxHubMessageReceiver} receives the CCTP V2 hook, validates the hook
recipient and mint recipient, executes hub calldata, and records failures as
stranded deposits. A grace-period sweep path lets the beneficiary recover
stranded funds after a hub-call failure.

\subsection{Circle Gateway for Hub-to-Hub USDC}

Circle Gateway is modeled as a fast USDC liquidity rail between \protocol{}
hubs. In the current design:

\begin{itemize}
  \item Gateway is USDC-only.
  \item The source signer is an EOA.
  \item ERC-1271 contract signing is modeled as a future mode but disabled
        until Circle support is live and allowlisted.
  \item The destination hook calls Gateway Minter, checks the exact USDC
        balance delta, records receipt state, and forwards USDC to the
        destination hub unless the hook itself is the destination.
\end{itemize}

\texttt{TelaranaGatewayHubHook} validates:

\begin{enumerate}
  \item route id,
  \item source and destination Gateway domains,
  \item destination Gateway Minter,
  \item destination USDC token,
  \item route caller authorization,
  \item nonzero amount and participants,
  \item request replay status,
  \item exact minted amount,
  \item action-specific spot-FX fields,
  \item undefined \texttt{hookData}, which is rejected until it has audited
        semantics.
\end{enumerate}

\subsection{Hyperlane for Intents and Approved Asset Routes}

Hyperlane is the intent and non-Gateway / non-CCTP route layer. The key lesson
from Velodrome's Superchain architecture is to avoid treating cross-chain as
``bridge everything everywhere.'' Instead, route user intent into the deepest
canonical market and settle accounting at the hub.

The Hyperlane flow is:

\[
\text{origin spoke} \rightarrow \text{typed intent} \rightarrow
\text{hub receiver validation} \rightarrow \text{registered hub action}.
\]

Warp Routes or issuer-specific routes may move non-USDC assets, but the hub
accepts them only after route, asset, oracle, and market approval.

\section{Uniswap v4 Spot-FX Execution}

\protocol{} prepares spot FX through Uniswap v4. The current hook work is not a
full final FX exchange; it is a risk-aware foundation:

\begin{itemize}
  \item oracle-anchored quotes,
  \item truncated observations,
  \item volatility-aware spread controls,
  \item reserve and rehypothecation accounting,
  \item future spot request events,
  \item future RFQ Pasillo quote corridor.
\end{itemize}

Uniswap v4 hooks are useful because pool lifecycle callbacks let protocol logic
run at initialization, liquidity modification, swap, and donation boundaries.
\protocol{} prioritizes the hook ideas most relevant to FX credit:

\begin{enumerate}
  \item truncated/geomean oracle concepts,
  \item dynamic volatility spread concepts,
  \item TWAMM / long-order concepts for larger FX execution,
  \item rehypothecation concepts that pair LP liquidity with Morpho supply.
\end{enumerate}

The hook surface is intentionally modular because future FX perpetuals may use
the same liquidity base. This paper does not specify or implement perpetuals;
it preserves compatibility by avoiding assumptions that would make the
liquidity unusable for a later trading engine.

\section{Ghost Mode}

Ghost Mode is the future private route family. It is not a cosmetic dark theme;
it is a slower, separate execution mode for users with a valid Bufi Wallet
KYC/KYB pass.

Ghost Mode requires:

\begin{enumerate}
  \item Bufi Wallet pass verification,
  \item a trusted Bufi router,
  \item commitment/nullifier accounting,
  \item separate privacy-capable hook or pool surfaces,
  \item the same hub risk engine used by public mode.
\end{enumerate}

The public pools remain permissionless. Ghost pools or routes can be
pass-gated. KYC examples that rely on \texttt{tx.origin} are not suitable for
production. In Uniswap v4, hooks must account for PoolManager and router call
semantics instead of assuming the externally owned account is the direct caller.

\section{Governance and Operations}

The protocol favors operational levers that are narrow and observable:

\begin{itemize}
  \item role-gated admin and operations actions,
  \item global pause for entry-side actions,
  \item per-pool live switches,
  \item deployment manifests,
  \item SDK ABI generation,
  \item indexer-ready events,
  \item stranded-deposit recovery,
  \item route metadata and event schemas.
\end{itemize}

The design avoids premature token governance. Governance, incentives, gauges,
or voting escrow systems can be considered after markets are safe, liquid, and
observable.

\section{Verification Status}

The current smart-contract protocol handoff scope is marked complete in
\texttt{docs/BUCKET\_ANALYSIS\_SMART\_CONTRACTS\_2026-05-15.md}. The observed
verification runs were:

\begin{center}
\begin{tabular}{lll}
\toprule
Command & Result & Notes \\
\midrule
\texttt{forge test --match-contract TelaranaGatewayHubHookTest -vvv} &
17 passed & Gateway edge cases. \\
\texttt{bun run contracts:test} & 147 passed, 1 skipped & Local unit and
invariant suite; skip is optional Tenderly manifest. \\
\texttt{bun run contracts:test:fork} & 161 passed, 1 skipped & Includes
mainnet-fork Morpho tests. \\
\texttt{bun run sdk:test} & 33 passed & SDK, ABI, route, Gateway helpers. \\
\texttt{bun run sdk:build} & Passed & TypeScript build. \\
\bottomrule
\end{tabular}
\end{center}

\section{Gateman Verification Discipline}

This work used the \texttt{gateman-analysis} skill, inspired by Professor
Robert Gateman's engineering-translated discipline:

\begin{quote}
Assume nothing, question everything, worship no one, and applaud humility.
\end{quote}

In this repository, the Gateman pass closed five concrete findings:

\begin{enumerate}
  \item imprecise Circle/Gateway rail language,
  \item undefined \texttt{hookData} semantics,
  \item mint-only Gateway requests carrying spot fields,
  \item same-domain Gateway routes,
  \item insufficient Gateway unhappy-path tests.
\end{enumerate}

The protocol should keep this posture. A whitepaper, a passing test suite, or a
famous dependency is not proof of safety. Evidence is proof.

\section{Limitations and Open Release Controls}

The following are not complete merely because the smart-contract handoff is
complete:

\begin{itemize}
  \item external audit,
  \item production deployment,
  \item production Circle Gateway configuration,
  \item Circle ERC-1271 contract-signing support,
  \item Hyperlane route deployments,
  \item route-specific ISM review,
  \item relayer monitoring,
  \item live oracle-feed verification on each target chain,
  \item incident response rehearsal,
  \item formal verification of critical invariants.
\end{itemize}

\section{Conclusion}

\protocol{} is a liquidity web for FX credit, not a bridge aggregator. Its
core contribution is the separation of transport from risk acceptance. CCTP,
Gateway, Hyperlane, issuer routes, Morpho, Pyth, RedStone, and Uniswap v4 each
solve different parts of the system. \protocol{} composes them around one
principle: the hub decides what is safe to accept, borrow against, swap through,
or route onward.

The result is an architecture that can start with conservative USDC entry and
the deepest available fiat-stablecoin hub markets, then expand toward
hub-to-hub USDC liquidity, non-USDC approved asset routes, private Ghost Mode
execution, RFQ Pasillo, and future trading frontends without relaxing the core
collateral-safety model.

\section*{Acknowledgements}

We give thanks to God for truth, discipline, and the obligation to steward risk
carefully.

We thank Professor Robert Gateman for the four-law posture that shaped the
repository's post-implementation audit discipline. We thank the authors and
maintainers of the \texttt{gateman-analysis} skill for translating that posture
into a practical engineering workflow.

We thank Tom\'as Cordero for the protocol direction, requirements, market
references, architecture pressure, and product framing. We thank the
OpenAI Codex/GPT engineering stack for the model-assisted implementation work.

We thank the builders whose public work informed this repository: Morpho,
Uniswap, Circle, Hyperlane, Velodrome, Pyth, RedStone, OpenZeppelin, Austin
Adams and the Uniswap hook-example contributors, BlackBera's privacy-hook
concept, J. Dub Park's public KYC hook example, and Constantino's
Uniswap-v4/Morpho rehypothecation hook work.

\appendix

\section{Model Contribution and Token Ledger}
\label{appendix:token-ledger}

This paper was drafted by GPT.5.5 Extra High as the main author label requested
by the maintainer, with Tom\'as Cordero as co-author and protocol owner.

Exact token usage for the model contribution is not available from repository
state, Git metadata, or the local shell runtime. The verified token ledger for
publication is therefore:

\begin{center}
\begin{tabular}{ll}
\toprule
Field & Value \\
\midrule
Model author label & GPT.5.5 Extra High \\
Human co-author & Tom\'as Cordero \\
Repository & \url{https://github.com/BuFi007/fx-telarana} \\
Protocol verification commit & \commithead{} \\
Verified model tokens used & \tokenledger{} \\
\bottomrule
\end{tabular}
\end{center}

The token field must be replaced with the value from the provider usage export
before publication. No estimated or invented token total is reported here.

\section{Repository Contribution Ledger}

\begin{center}
\begin{tabular}{llp{0.52\linewidth}}
\toprule
Commit & Type & Contribution \\
\midrule
\texttt{b19cdab} & Feature & Prepared Circle Gateway hub-liquidity SDK
config, interfaces, typed data, and docs. \\
\texttt{297f622} & Feature & Added \texttt{TelaranaGatewayHubHook}, tests,
and ABI export. \\
\texttt{8de5959} & Verification & Closed smart-contract bucket analysis,
Gateway edge hardening, ABI regeneration, and Gateman report. \\
\bottomrule
\end{tabular}
\end{center}

\section{Reference Repositories and Documents}

\begin{thebibliography}{99}

\bibitem{telarana}
BuFi007, \textit{fx-telarana repository}. \url{https://github.com/BuFi007/fx-telarana}

\bibitem{morpho-blue}
Morpho Labs, \textit{Morpho Blue}. \url{https://github.com/morpho-org/morpho-blue}

\bibitem{uniswap-v4-core}
Uniswap Labs, \textit{Uniswap v4 Core}. \url{https://github.com/Uniswap/v4-core}

\bibitem{uniswap-v4-periphery}
Uniswap Labs, \textit{Uniswap v4 Periphery}. \url{https://github.com/Uniswap/v4-periphery}

\bibitem{circle-gateway}
Circle, \textit{EVM Gateway Contracts}. \url{https://github.com/circlefin/evm-gateway-contracts}

\bibitem{circle-cctp}
Circle, \textit{CCTP EVM Contracts}. \url{https://github.com/circlefin/evm-cctp-contracts}

\bibitem{hyperlane}
Hyperlane, \textit{Hyperlane Monorepo}. \url{https://github.com/hyperlane-xyz/hyperlane-monorepo}

\bibitem{hyperlane-docs}
Hyperlane, \textit{Documentation}. \url{https://docs.hyperlane.xyz/}

\bibitem{velodrome-superchain}
Velodrome Finance, \textit{Superchain Contracts Specification}. \url{https://github.com/velodrome-finance/superchain-contracts}

\bibitem{pyth}
Pyth Network, \textit{Pyth Crosschain}. \url{https://github.com/pyth-network/pyth-crosschain}

\bibitem{redstone}
RedStone Finance, \textit{RedStone Oracles Monorepo}. \url{https://github.com/redstone-finance/redstone-oracles-monorepo}

\bibitem{openzeppelin}
OpenZeppelin, \textit{OpenZeppelin Contracts}. \url{https://github.com/OpenZeppelin/openzeppelin-contracts}

\bibitem{flashifi}
Constantino, \textit{flashifi-uniswap-v4-rehypothecation-hook}. \url{https://github.com/Constantino/flashifi-uniswap-v4-rehypothecation-hook}

\bibitem{privacy-hook}
BlackBera, \textit{privacy-hook-univ4}. \url{https://github.com/blackbera/privacy-hook-univ4}

\bibitem{kyc-hook}
J. Dub Park, \textit{Uniswap-Hooks / KYCHook}. \url{https://github.com/jdubpark/Uniswap-Hooks}

\bibitem{gateman-skill}
\texttt{gateman-analysis}, \textit{Post-implementation code audit skill}. Local path:
\texttt{/Users/criptopoeta/.codex/skills/gateman-analysis/SKILL.md}

\bibitem{gateman-humans}
Humans of UBC, \textit{Assume Nothing: Robert Gateman, Economics Professor}. \url{https://www.facebook.com/humansofubc/posts/assume-nothing-robert-gateman-economics-professor/250788225628894/}

\bibitem{gateman-reddit}
r/UBC, \textit{Motivational Gateman quotes}. \url{https://www.reddit.com/r/UBC/comments/6lb7al/motivational_gateman_quotes/}

\bibitem{gateman-blog}
Kiely, \textit{Gateman's Econ 101}. UBC Blogs. \url{https://blogs.ubc.ca/kiely/2011/11/27/gatemans-econ-101/}

\end{thebibliography}

\end{document}
